\section{Removing threshold-weight magnitudes from the teacher parameters}
\label{sec:critical-weight-independent-threshold}

A Boolean threshold gate may be specified using arbitrarily large or
arbitrarily precise real coefficients. These numerical parameters are
unnecessary for the learning bound. The gate's Boolean input set admits
an equivalent separator with polynomially many bits, and constant-depth
majority arithmetic evaluates that separator using polynomially bounded
gate coefficients. This is an existence result; it does not require an
algorithm that receives or computes with the original real coefficients.

\begin{lemma}[Small-bit integer separators]
\label{lem:critical-small-bit-separator}
Let $S\subseteq\{-1,1\}^k$, where $k\ge0$, and suppose
$g:S\to\{-1,1\}$ agrees with an affine real threshold, with either fixed
choice of sign at zero. There are integers $b_0,\ldots,b_k$ such that
\[
 g(u)\left(b_0+\sum_{j=1}^k b_j u_j\right)\ge1
 \quad(u\in S),\qquad
 |b_j|\le (k+1)^{(k+1)/2}.
\]
In particular, each coefficient has
$O((k+1)\log(k+2))$ bits, with a universal constant.
The assertion applies to the full Boolean cube or to only the states
realized by earlier gates.
\end{lemma}

\begin{proof}
The empty set is immediate, so assume $S$ is nonempty. If any original
scores vanish, change the bias by a sufficiently small quantity whose
sign is the specified sign at zero. Finiteness of $S$ ensures that this
does not change the sign of any nonzero score. All desired inequalities
then become strict. Scale the coefficients so that every signed score
is at least one.

Put $m=k+1$ and form the matrix $A$ whose row for $u\in S$ is
$g(u)(1,u_1,\ldots,u_k)$. Thus $A$ has entries in $\{-1,1\}$ and
$Aw\ge\boldsymbol1$ is feasible. Let $r=\operatorname{rank}A$, so
$1\le r\le m$, and retain $r$ linearly independent columns to obtain
$A_J$. These columns span the column space of $A$, so
$A_Jv\ge\boldsymbol1$ remains feasible.

We verify directly that this latter system has a feasible point with
$r$ linearly independent tight inequalities. Start at any feasible
point. If its tight rows have rank less than $r$, choose a nonzero
direction orthogonal to all of them. Since $A_J$ has column rank $r$,
some row has a nonzero inner product with the direction. Reverse the
direction if necessary so that at least one row has negative inner
product. Every such row is currently slack. Move until the first of
these finitely many rows becomes tight: the permitted distance is the
positive minimum of its slack divided by the negative directional
derivative. All old tight rows stay tight, feasibility is preserved,
and a new tight row is independent of them because it is not
orthogonal to the chosen direction. After at most $r$ iterations the
tight rows have rank $r$.

Choose $r$ independent tight rows, forming a nonsingular integer
matrix $M\in\{-1,1\}^{r\times r}$. At the resulting feasible point,
$Mv=\boldsymbol1$. Cramer's rule gives
\[
 v_j=\frac{\det M_j}{\det M},
\]
where $M_j$ replaces column $j$ by $\boldsymbol1$. Multiply the
feasible coefficient vector by $|\det M|$, and put zero in the omitted
coordinates. The retained coefficients are the integers
$\operatorname{sign}(\det M)\det M_j$. Their signed scores are at
least $|\det M|\ge1$. All entries of $M_j$ still have absolute value
one. Hadamard's determinant bound therefore gives
\[
 |\det M_j|\le r^{r/2}\le m^{m/2},
\]
as required. This also covers $k=0$, where $r=m=1$.
\end{proof}

\begin{lemma}[A weight-independent threshold compiler]
\label{lem:critical-weight-independent-threshold-compiler}
There is a universal integer $a$ with the following property. Suppose
a Boolean target on $\{-1,1\}^d$ has a layered real threshold circuit
of depth $h\ge1$ and width at most $C\ge1$. Threshold gates may have
arbitrary real coefficients and biases, and may use either fixed tie
convention. Put $K=d+C+h+2$.
The target has a layered Boolean threshold implementation with depth
at most $ah$, width at most $K^a$, and Euclidean norm of every gate
coefficient vector at most $K^a$. Every gate has signed score at least
one on its realized Boolean states. Constant wires and layering
overhead are included. The bounds are independent of all original
coefficient magnitudes and precisions.
\end{lemma}

\begin{proof}
Apply Lemma~\ref{lem:critical-small-bit-separator} separately to every
gate, either on its full Boolean input cube or on its realized input
states. The maximum fan-in is at most $\max\{d,C\}$, after combining
duplicate wires; biases are included in the constant coefficient.
Thus every gate has an equivalent integer description with
polynomially many bits in $K$. Choosing these descriptions is
nonuniform and requires no representation of the original real
numbers.

Use the standard constant-depth majority-arithmetic theorem: iterated
addition and comparison of polynomially many polynomial-bit integers
have polynomial-size, constant-depth circuits of unweighted majority
and Boolean gates, with universal depth and polynomial degree. This
is part of the usual arithmetic threshold-circuit toolkit; the
stronger division and iterated-multiplication theorem is proved in
Hesse--Allender--Barrington~\cite{hesse2002}. To apply it here, hardcode
the integer coefficients in binary. Each input bit selects either
zero or the corresponding fixed summand. Sum positive and negative
parts separately and compare them, including the constant term. All
intermediate integers have polynomially many bits. This construction
uses the binary coefficients as input constants to the arithmetic
circuit; it does not put exponentially large bit-position weights
on the new threshold gates.

Replace each original gate by this constant-depth majority circuit.
All gates from one original level are evaluated in parallel. The
total depth is $O(h)$, and the total number of gates and wires is a
universal polynomial in $K$. Insert identity gates where necessary
to obtain a layered circuit and propagate a constant wire. This
costs only a polynomial in the size and depth.

Finally, every unweighted majority, AND, OR, or NOT gate has an
integer threshold realization with coefficients polynomially bounded
in its fan-in. Put its threshold halfway between consecutive allowed
integer sums and multiply the score by two if needed. On all Boolean
inputs its score has absolute value at least one and the prescribed
sign. Its bias is supplied by the constant wire, so its coefficient
Euclidean norm is polynomial in the total circuit size. Increasing
the universal exponent $a$ proves the stated simultaneous bounds.
\end{proof}

This replacement acts on Boolean states. It does not remove a weight
bound from a real-valued tanh teacher: preserving intermediate real
activations is a different approximation problem.
